% =============================================================================
% AegisMAS Report — Section: Attack Design (First Draft)
% Author: Crispin (Group Leader + Attack Dev)
% For: Overleaf (IEEE IEEEtran format)
% =============================================================================
% Paste this into your Overleaf document in the appropriate section.
% Adjust the section number to match your report outline.
% =============================================================================

\section{Attack Design}
\label{sec:attacks}

This section describes the threat model and the four attack conditions
used to evaluate the resilience of the multi-agent system.

% -----------------------------------------------------------------------------
\subsection{Threat Model}
\label{subsec:threat-model}

We consider a multi-agent supply-chain logistics system in which
a set of specialised agents (warehouse, transport, and route agents)
report observations to a central coordinator.
The coordinator aggregates these \emph{claims} and builds a delivery plan.
A plan checker then validates the plan against ground truth
(i.e., whether the plan would actually succeed if executed).

Our threat model makes the following assumptions:

\begin{enumerate}
    \item \textbf{Claim-level manipulation.}
        The adversary can modify the claims that agents produce
        \emph{after} the agent reports but \emph{before} the coordinator
        receives them.  This is the ``wrapper'' pattern: the agent code
        itself is never altered, so the attack is external to the agent's
        logic.

    \item \textbf{Knowledge of ground truth.}
        The adversary has access to the true mission parameters
        (warehouse stock, truck capacity, route times) and can craft
        lies that are plausible yet harmful---e.g., inflating the stock
        of a warehouse that is actually empty.

    \item \textbf{Goal of the adversary.}
        The adversary's goal is to cause the coordinator to adopt
        a plan that \emph{appears} valid based on the (corrupted) claims
        but \emph{fails} when checked against reality.
        A successful attack produces a plan that the plan checker
        rejects.

    \item \textbf{No defence awareness (for now).}
        In the base attack conditions, the adversary does not adapt
        to the defence.  This is a standard first-pass assumption;
        adaptive attacks are noted as future work (Section~\ref{sec:future}).
\end{enumerate}

% -----------------------------------------------------------------------------
\subsection{Attack Conditions}
\label{subsec:attack-conditions}

We define four conditions, labelled C0--C3, that span a spectrum from
no interference to coordinated multi-vector attacks.
Table~\ref{tab:attack-summary} summarises them.

\begin{table}[htbp]
\centering
\caption{Summary of attack conditions.}
\label{tab:attack-summary}
\begin{tabular}{c l l}
\hline
\textbf{ID} & \textbf{Name} & \textbf{What is corrupted} \\
\hline
C0 & No attack (control)     & Nothing \\
C1 & Byzantine agent         & One agent's claims \\
C2 & Poisoned tool           & One tool's observations \\
C3 & Coordinated             & Agent + tool reinforce same lie \\
\hline
\end{tabular}
\end{table}

% --- C0 ---
\subsubsection{C0 --- No Attack (Control Baseline)}
\label{subsubsec:c0}

Under C0 all agents report truthfully and no claims are modified.
This condition serves two purposes:
(a)~it measures the \emph{baseline task success rate}
    (TSR) of the system when no adversary is present, and
(b)~it quantifies the \emph{false-positive rate} (FPR) of each defence---
    any claims flagged as suspicious under C0 are false alarms.

\subsubsection{C1 --- Byzantine Agent}
\label{subsubsec:c1}

In C1 a single agent is selected as the \emph{Byzantine} agent.
Specifically, one warehouse agent is compromised:
it reports inflated stock levels so that the coordinator
believes that warehouse can fulfil the order when, in reality,
it cannot.

The attack proceeds as follows:
\begin{enumerate}
    \item Select the \emph{weakest} warehouse---the one whose true stock
          is furthest below the mission's target requirements.
    \item Replace its \texttt{medical\_kits} and \texttt{tents} claims
          with values that exceed the target by a comfortable margin
          (e.g., target $+ 50$ kits, target $+ 30$ tents).
    \item Leave all other claims (transport, route, other warehouses)
          unchanged.
\end{enumerate}

The coordinator, seeing apparently sufficient stock at the lying
warehouse, may choose it for the plan.  When the plan checker
validates against the true (insufficient) stock, the plan fails.

This attack models the classic Byzantine fault scenario~\cite{lamport1982byzantine}
in which a single participant sends incorrect information
to mislead the group's decision.

\subsubsection{C2 --- Poisoned Tool}
\label{subsubsec:c2}

In C2 no agent is dishonest.  Instead, a \emph{tool}---the route
information source---is poisoned.  The route agent queries the tool
honestly but receives manipulated data, which it then relays
in good faith to the coordinator.

The attack modifies route claims as follows:
\begin{enumerate}
    \item Select the route to the warehouse with the \emph{lowest} true
          stock (the worst option for the coordinator to choose).
    \item Reduce the reported \texttt{time\_hours} so the route appears
          fast enough to meet the deadline, even if the real travel
          time would exceed it.
    \item Optionally reduce \texttt{distance\_km} for consistency
          (a short route should also be fast).
\end{enumerate}

The coordinator now believes that a warehouse which is either
understocked or too far away is in fact reachable in time, and may
route the delivery through it.  The honest agent is an unwitting
accomplice---it reported exactly what its (poisoned) tool told it.

This condition reflects the growing concern around
tool-poisoning and indirect prompt injection
in LLM-based agent systems~\cite{greshake2023youve},
adapted here to a non-LLM scripted setting.

\subsubsection{C3 --- Coordinated Attack (Agent + Tool)}
\label{subsubsec:c3}

C3 combines C1 and C2: a Byzantine agent \emph{and} a poisoned tool
reinforce the same false narrative.  Both vectors target the same
warehouse:

\begin{enumerate}
    \item The warehouse agent lies about stock levels (as in C1),
          making the target warehouse appear well-stocked.
    \item The route tool lies about travel time (as in C2),
          making the route to that warehouse appear fast.
\end{enumerate}

Because both the stock report and the route report agree, the
coordinator receives a \emph{consistent} but entirely false picture.
This makes the lie harder to detect:
a defence that cross-checks stock against route information
would find no internal contradiction.

C3 tests whether defences that perform well against single-vector
attacks (C1 or C2 alone) degrade when confronted with coordinated
deception.

% -----------------------------------------------------------------------------
\subsection{Attack Implementation}
\label{subsec:attack-impl}

All attacks are implemented as pure functions with the signature:

\begin{verbatim}
attack_fn(mission, claims) -> modified_claims
\end{verbatim}

\noindent
where \texttt{mission} provides ground truth (so the adversary can
craft targeted lies) and \texttt{claims} is the list of agent reports.
The function returns a new list of claims with selected entries
modified.  Unmodified claims are passed through unchanged.

This wrapper design has three advantages:
\begin{itemize}
    \item \textbf{Separation of concerns:}
        agent code never contains attack logic; attacks are external
        interceptors.
    \item \textbf{Composability:}
        C3 is built by applying C1's manipulation followed by C2's,
        rather than requiring a separate monolithic function.
    \item \textbf{Reproducibility:}
        because the attack function is deterministic given the mission
        and claims, the same seed always produces the same corrupted
        scenario.
\end{itemize}

% -----------------------------------------------------------------------------
\subsection{Attack Targets and Selection Strategy}
\label{subsec:target-selection}

A key design decision is \emph{which} agent or tool to compromise.
Random selection would add noise to the results: a randomly chosen
warehouse that happens to already have sufficient stock cannot
cause a plan failure by inflating its numbers further.

We therefore use a \textbf{worst-case targeting} strategy:
the attack always selects the target that maximises the chance of
misleading the coordinator.  For C1, this means the warehouse with
the lowest true stock; for C2, the route with the longest true
travel time.  This gives an upper bound on attack effectiveness
and ensures the experiment measures defence capability against
a competent (not random) adversary.

% -----------------------------------------------------------------------------
\subsection{Summary}

The four conditions form a $2 \times 2$ design over two attack
vectors (agent compromise and tool poisoning), as shown in
Table~\ref{tab:attack-2x2}.

\begin{table}[htbp]
\centering
\caption{Attack conditions as a $2 \times 2$ design.}
\label{tab:attack-2x2}
\begin{tabular}{l c c}
\hline
 & \textbf{Tool honest} & \textbf{Tool poisoned} \\
\hline
\textbf{Agent honest}      & C0 (control) & C2 \\
\textbf{Agent Byzantine}   & C1           & C3 \\
\hline
\end{tabular}
\end{table}

\noindent
This factorial structure allows the experiment to measure the
\emph{individual} effect of each vector (C1 vs.\ C0, C2 vs.\ C0)
as well as their \emph{interaction} (C3 vs.\ the sum of C1 and C2
effects), which directly addresses the ``coordinated'' aspect of
the research question.
